\begin{center}
\hypertarget{native-037-article019}{}\pdfbookmark[2]{19.}{article-019}19.
\end{center}
Solutio problematis, ex anomalia vera deducere tempus, multoque adhuc magis solutio problematis inversi, magnopere abbreviari potest per tabulam auxiliarem, qualis in pluribus libris astronomicis reperitur. Longe vero commodissima est tabula \textsc{Barker}iana, quae etiam operi egregio cel. \textsc{Olbers} (\emph{Abhandlung über die leichteste und bequemste Methode die Bahn eines Cometen zu berechnen}, Weimar 1797.) annexa est. Continet ea pro omnibus anomaliis veris a 0 usque ad $180^0$ per singula 5 minuta valorem expressionis $75\operatorname{tang}\frac{1}{2}v+25\operatorname{tang}\frac{1}{2}v^3$ sub nomine \emph{motus medii}. Si itaque tempus desideratur anomaliae verae $v$ respondens, dividere oportebit motum medium e tabula argumento $v$ excerptum per $\frac{150k}{p^{\frac{3}{2}}}$, quae quantitas \emph{motus medius diurnus} dicitur; contra si e tempore anomalia vera computanda est, illud in diebus expressum per $\frac{150k}{p^{\frac{3}{2}}}$ multiplicabitur, ut motus medius prodeat, quo anomaliam respondentem e tabula sumere licebit. Ceterum manifesto valori negativo ipsius $v$ motus medius tempusque idem sed negative sumtum respondet: eadem igitur tabula anomaliis negativis et positivis perinde inservit. Si pro $p$ distantia in perihelio $\frac{1}{2}p=q$ uti malumus, motus medius diurnus exprimitur per $\frac{k\BareRoot{2812{,}5}}{q^{\frac{3}{2}}}$, ubi factor constans $k\BareRoot{2812{,}5}$ fit $=0{,}912\,2790\,61$, ipsiusque logarithmus $9{,}960\,1277\,069$. — Inventa anomalia $v$ radius vector determinabitur per formulam iam supra traditam $r=\frac{q}{\cos\frac{1}{2}v^2}\cdot$

\begin{center}
\hypertarget{native-037-article020}{}\pdfbookmark[2]{20.}{article-020}20.
\end{center}
Per differentiationem aequationis $\operatorname{tang}\frac{1}{2}v+\frac{1}{3}\operatorname{tang}\frac{1}{2}v^3=2t k p^{-\frac{3}{2}}$, si omnes quantitates $v,t,p$ ceu variabiles tractantur, prodit
\[
\frac{\mathrm{d}v}{2\cos\frac{1}{2}v^4}=2k p^{-\frac{3}{2}}\,\mathrm{d}t-3t k p^{-\frac{5}{2}}\,\mathrm{d}p
\]
sive
\[
\mathrm{d}v=\frac{k\BareRoot{p}}{r r}\,\mathrm{d}t-\frac{3t k}{2r r\BareRoot{p}}\,\mathrm{d}p.
\]
\SourceLeaf{038}Si variationes anomaliae $v$ in secundis expressae desiderantur, etiam ambae partes ipsius $\mathrm{d}v$ hoc modo exprimendae sunt, i. e. pro $k$ valorem in art. 6 traditum $3548^{\prime\prime}{,}188$ accipere oportet. Quodsi insuper pro $p$ introducatur $\frac{1}{2}p=q$, formula ita se habebit
\[
\mathrm{d}v=\frac{k\BareRoot{2q}}{r r}\,\mathrm{d}t-\frac{3k t}{r r\BareRoot{2q}}\,\mathrm{d}q,
\]
ubi logarithmi constantes adhibendi sunt $\log k\BareRoot{2}=3{,}700\,5215\,724$,\\ $\log 3k\BareRoot{\frac{1}{2}}=3{,}876\,6128\,315$.

Porro differentiatio aequationis $r=\frac{p}{2\cos\frac{1}{2}v^2}$ suppeditat
\[
\frac{\mathrm{d}r}{r}=\frac{\mathrm{d}p}{p}+\operatorname{tang}\frac{1}{2}v\,\mathrm{d}v
\]
sive exprimendo $\mathrm{d}v$ per $\mathrm{d}t$ et $\mathrm{d}p$
\[
\frac{\mathrm{d}r}{r}=\left(\frac{1}{p}-\frac{3k t\operatorname{tang}\frac{1}{2}v}{2r r\BareRoot{p}}\right)\mathrm{d}p+\frac{k\BareRoot{p}\operatorname{tang}\frac{1}{2}v}{r r}\,\mathrm{d}t.
\]
Coëfficiens ipsius $\mathrm{d}p$, substituendo pro $t$ valorem suum per $v$, transit in
\[
\begin{aligned}
&\frac{1}{p}-\frac{3p\operatorname{tang}\frac{1}{2}v^2}{4r r}-\frac{p\operatorname{tang}\frac{1}{2}v^4}{4r r}\\
&\quad=\frac{1}{r}\left(\frac{1}{2}+\frac{1}{2}\operatorname{tang}\frac{1}{2}v^2-\frac{3}{2}\sin\frac{1}{2}v^2-\frac{1}{2}\sin\frac{1}{2}v^2\operatorname{tang}\frac{1}{2}v^2\right)\\
&\quad=\frac{\cos v}{2r};
\end{aligned}
\]
coëfficiens ipsius $\mathrm{d}t$ autem fit $=\frac{k\sin v}{r\BareRoot{p}}\cdot$ Hinc prodit $\mathrm{d}r=\frac{1}{2}\cos v\,\mathrm{d}p+\frac{k\sin v}{\BareRoot{p}}\,\mathrm{d}t$ sive, introducendo $q$ pro $p$,
\[
\mathrm{d}r=\cos v\,\mathrm{d}q+\frac{k\sin v}{\BareRoot{2q}}\,\mathrm{d}t.
\]
Logarithmus constans hic adhibendus est $\log k\BareRoot{\frac{1}{2}}=8{,}085\,0664\,436$.
